\section{Represent the outgoing branch as a square-integrable eigenvector}
\label{sec:complex-scaling}
% BEGIN CURATED SUBJECT INDEX
\index{operator!closed}
\index{spectrum!continuous}
\index{spectrum!point}
\index{resolvent}
\index{Riesz projection}
\index{boundary condition!incoming and outgoing}
% END CURATED SUBJECT INDEX

\calculationroute
  {The continued pole condition, closed operators, analytic families, and
  the distinction between point and continuous spectrum.}
  {Choose an admissible coordinate deformation that makes the outgoing branch
  square integrable while rotating, rather than discarding, the continuum.}
  {Identify an angle-stable eigenvalue and connect its Riesz projection to the
  same continued resolvent pole.}

\subsection{Analytic dilation}
\label{subsec:analytic-dilation}
% BEGIN CURATED SUBJECT INDEX
\index{scattering!Coulomb}
\index{complex scaling}
\index{analytic dilation}
\index{complex scaling!exterior}
% END CURATED SUBJECT INDEX

Let the configuration-space dimension be $d$.  For a real dilation parameter
$s$, define the unitary operator
\begin{equation}
  \left[U(s)\psi\right](\bm x)
  =\rme^{ds/2}\psi(\rme^s\bm x) .
  \label{eq:real-dilation}
\end{equation}
Complex scaling analytically continues $s$ to $s=\rmi\theta$, with a real
angle $\theta$.  Formally,
\begin{equation}
  \left[U(\rmi\theta)\psi\right](\bm x)
  =\rme^{\rmi d\theta/2}
  \psi(\rme^{\rmi\theta}\bm x) .
  \label{eq:complex-dilation}
\end{equation}
For $\theta\ne0$, this transformation is not unitary on the original
$L^2$ space and is generally unbounded.  It should be defined first on a dense
set of analytic vectors.  This is why it is safer to speak of an analytic
family of deformed operators than of an everywhere-defined similarity
transformation.

For the Hamiltonian
\begin{equation}
  H=T+V(\bm x),
  \qquad
  T=-\frac{\hbar^2}{2m}\nabla^2,
  \label{eq:H-TV}
\end{equation}
the complex-scaled operator is
\begin{equation}
  H_{\theta}
  =U(\rmi\theta)HU(-\rmi\theta)
  =\rme^{-2\rmi\theta}T
  +V(\rme^{\rmi\theta}\bm x) .
  \label{eq:H-theta}
\end{equation}
Equation~\eqref{eq:H-theta} requires the potential to admit analytic
continuation in the sector swept out by the deformation.  Coulomb singularities,
multiple coordinates, and nonanalytic potentials require refined forms such
as exterior complex scaling or smooth complex deformation
\cite{Simon1979ECS,Reinhardt:1982,Moiseyev:2011}.

\subsection{Why an outgoing wave becomes square integrable}
\label{subsec:outgoing-L2}
% BEGIN CURATED SUBJECT INDEX
\index{boundary condition!incoming and outgoing}
\index{Gamow--Siegert state}
% END CURATED SUBJECT INDEX

Write the resonance momentum as
\begin{equation}
  \vsub{k}{R}=|\vsub{k}{R}|\rme^{-\rmi\varphi},
  \qquad 0<\varphi<\frac{\pi}{4} .
  \label{eq:k-polar}
\end{equation}
On the positive half-line, the scaled outgoing wave is
\begin{align}
  \rme^{\rmi \vsub{k}{R}r\rme^{\rmi\theta}}
  &=\rme^{
  \rmi|\vsub{k}{R}|r\cos(\theta-\varphi)
  }
  \notag\\
  &\quad\times
  \rme^{-|\vsub{k}{R}|r
  \sin(\theta-\varphi)} .
  \label{eq:scaled-outgoing}
\end{align}
It decays for $\theta>\varphi$.  The same condition regularizes the outgoing
wave on the negative half-line when the contour is rotated consistently.
Thus the Siegert state becomes an $L^2$ eigenfunction of $H_\theta$ once the
rotated continuum has passed below the pole.  The divergent unscaled wave and
the normalizable scaled wave are analytic continuations of one solution, not
different physical states.

\subsection{The Aguilar--Balslev--Combes structure}
\label{subsec:abc}
% BEGIN CURATED SUBJECT INDEX
\index{operator!adjoint and self-adjoint}
\index{unbounded operator}
\index{operator!closed}
\index{spectrum!continuous}
\index{spectrum!essential}
\index{spectrum!point}
\index{resolvent}
\index{Riesz projection}
\index{pseudospectrum}
\index{compact operator}
% END CURATED SUBJECT INDEX

The Aguilar--Balslev--Combes theorem is sometimes quoted only as the rule that
the continuum rotates.  Its real content is a statement about an analytic
family of closed operators and the meromorphic continuation of resolvent
matrix elements
\cite{Aguilar:1971ve,Balslev:1971vb,Simon1972BalslevCombes,Simon:73}.
We state a one-threshold version before explaining its proof.

Let $H_0=-\hbar^2\nabla^2/(2m)$ on $L^2(\real^d)$ and assume that $V$ is
$H_0$-compact, or satisfies a suitable relative-form compactness condition.
Assume further that
\begin{equation}
  s\longmapsto V_s:=U(s)VU(-s)
  \label{eq:dilation-analytic-V}
\end{equation}
has an operator- or form-valued analytic continuation from real $s$ to a
strip $\lvert\im s\rvert<\theta_{\max}$.  Such a potential is called
\textit{dilation analytic}.  The deformed Hamiltonian
\begin{equation}
  H(s)=\rme^{-2s}H_0+V_s
  \label{eq:analytic-Hs}
\end{equation}
then forms an analytic family, of type A when a common operator domain can be
used and of type B when the construction is made through quadratic forms.
For $s=\rmi\theta$ with $0<\theta<\theta_{\max}$, write
$H_\theta:=H(\rmi\theta)$.

Under these hypotheses, the theorem gives the following spectral picture:
\begin{enumerate}
  \item $H_\theta$ is a closed non-self-adjoint operator.  Its essential
  spectrum is
  \begin{equation}
    \essspectrum(H_\theta)
    =\rme^{-2\rmi\theta}[0,\infty) .
    \label{eq:abc-essential-spectrum}
  \end{equation}

  \item Bound-state eigenvalues remain on the real axis and are independent of
  $\theta$.  Their algebraic multiplicities are unchanged for an admissible
  deformation that does not encounter another threshold or singularity.

  \item A discrete nonreal eigenvalue in the wedge
  $-2\theta<\arg z<0$ is independent of $\theta$ while it remains isolated
  from the rotated essential spectrum.  It coincides with a pole of the
  continued physical resolvent and is therefore a resonance.

  \item For many-body or coupled-channel problems, each continuum ray starting
  at a threshold $E_a$ rotates as
  \begin{equation}
    E_a+[0,\infty)
    \longmapsto
    E_a+\rme^{-2\rmi\theta}[0,\infty) .
    \label{eq:continuum-rotation}
  \end{equation}

  Cluster thresholds and channel branch points therefore remain as anchors of
  the spectral geometry.
\end{enumerate}
For a threshold at zero, a pole with energy $\resonantE$ is exposed when
\begin{equation}
  -2\theta<\arg \resonantE<0 .
  \label{eq:exposure-condition}
\end{equation}
The bridge to scattering theory is obtained on a dense set $\mathcal A$ of
analytic vectors.  For real $s$ unitarity of $U(s)$ gives
\begin{equation}
  \bra{\phi}(z-H)^{-1}\ket{\psi}
  =\bra{U(s)\phi}(z-H(s))^{-1}U(s)\ket{\psi},
  \qquad \text{$\phi$, $\psi\in\mathcal A$} .
  \label{eq:abc-resolvent-identity-main}
\end{equation}
Both sides are analytic in a common initial domain.  Continuing
Eq.~\eqref{eq:abc-resolvent-identity-main} from real $s$ to
$s=\rmi\theta$ continues the physical matrix element through the original
continuous spectrum into the wedge uncovered by the rotated ray.  Analytic
Fredholm theory makes the right-hand side meromorphic there.  Its poles are
the isolated eigenvalues of $H_\theta$, with finite-rank Riesz projections.
This is stronger than a graphical rotation of a cut: it identifies the
deformed eigenvalue with a pole of the undeformed scattering problem.

The theorem should not be paraphrased as a general completeness theorem for
arbitrary Gamow states.  It describes the spectrum and resolvent continuation
under specific analyticity and compactness assumptions.  A resonance
expansion still contains the rotated continuum, and completeness of a chosen
eigenfunction expansion is an additional statement.  The Berggren relation
in Section~\ref{subsec:berggren-proof} is one such statement for radial
short-range problems.

The phrase ``$\theta$ independent'' has a precise but limited meaning.  In the
exact operator problem, an isolated exposed pole is invariant under admissible
changes of $\theta$.  In a truncated basis, every eigenvalue has residual
angle dependence.  A numerical pole is identified by a stationary trajectory
or a plateau under simultaneous variation of $\theta$, basis size, nonlinear
basis scales, and arithmetic precision.  The pseudospectrum provides a more
honest error diagnostic than eigenvalue motion alone for strongly
non-normal matrices
\cite{Davies:2002,Davies:07,Trefethen1999}.

A detailed proof architecture, including the essential-spectrum and
Fredholm steps, is given in Sec.~\ref{app:abc-sketch}.  That discussion
also explains why global complex scaling does not immediately provide an ABC
theorem for black-hole potentials written in tortoise coordinates.

\subsection{Exterior and smooth complex scaling}
\label{subsec:exterior-scaling}
% BEGIN CURATED SUBJECT INDEX
\index{spectrum!point}
\index{complex scaling}
\index{complex scaling!smooth}
\index{complex absorbing potential}
% END CURATED SUBJECT INDEX

Global scaling evaluates $V$ at complex coordinates everywhere.  If the
interaction has singularities or if one wishes to preserve an inner physical
region, choose a radius $R_0>0$ and deform only outside it.  A piecewise-linear
exterior contour may be written
\begin{equation}
  z(r)=
  \begin{cases}
    r, & 0\le r\le R_0,\\
    R_0+(r-R_0)\rme^{\rmi\theta}, & r>R_0.
  \end{cases}
  \label{eq:ecs-contour}
\end{equation}
A smooth contour avoids a derivative discontinuity at $R_0$.  Exterior
complex scaling is especially useful when the potential is analytic only in
the asymptotic region or when an inner wave function must retain a direct
probabilistic interpretation.  Perfectly matched layers and complex absorbing
potentials are closely related numerical devices, but their finite-parameter
eigenvalues must be extrapolated to distinguish physical poles from absorber
modes.

\subsection{Biorthogonality and the complex-symmetric case}
\label{subsec:biorthogonality}
% BEGIN CURATED SUBJECT INDEX
\index{biorthogonality}
\index{complex-symmetric operator}
\index{Berggren completeness relation}
\index{exceptional point}
\index{self-orthogonality}
% END CURATED SUBJECT INDEX

Let $H_\theta\ket{\Psi_n^{\mathrm R}}=E_n\ket{\Psi_n^{\mathrm R}}$ and
$\bra{\Psi_n^{\mathrm L}}H_\theta=E_n\bra{\Psi_n^{\mathrm L}}$.  Away from
degeneracies, left and right eigenvectors may be normalized by
\begin{equation}
  \braket{\Psi_m^{\mathrm L}|\Psi_n^{\mathrm R}}
  =\delta_{mn} .
  \label{eq:biorthonormality}
\end{equation}
If a real potential and basis make the matrix representation complex
symmetric, $H_\theta^{\mathsf T}=H_\theta$, the natural bilinear form is the
$c$ product
\begin{equation}
  \cprod{\phi}{\psi}
  =\int\rmd x\,\phi(x)\psi(x) ,
  \label{eq:c-product}
\end{equation}
without complex conjugation of the first factor.  The $c$ product is not a
positive norm.  It is a bilinear pairing that implements left--right
biorthogonality in this representation.  Confusing it with a probability norm
is a common source of paradoxes concerning complex expectation values
\cite{Berggren:1996,Moiseyev:2011}.

At an exceptional point, Eq.~\eqref{eq:biorthonormality} cannot be maintained
with bounded normalization because the eigenvector becomes self-orthogonal.
The divergent normalization and the vanishing bilinear norm are two sides of
the same Jordan singularity.  We return to this point in
Section~\ref{sec:exceptional-points}.

\subsection{Spectral representation of the complex-scaled resolvent}
\label{subsec:csm-resolvent}
% BEGIN CURATED SUBJECT INDEX
\index{resolvent}
% END CURATED SUBJECT INDEX

In a formal discrete-plus-continuum notation, the deformed resolvent has the
expansion
\begin{align}
  \Resolvent_\theta(z)
  &=(z-H_\theta)^{-1}
  \notag\\
  &=\sum_{b}\frac{\ket{\Psi_b^{\mathrm R}}
  \bra{\Psi_b^{\mathrm L}}}{z-E_b}
  +\sum_{r}\frac{\ket{\Psi_r^{\mathrm R}}
  \bra{\Psi_r^{\mathrm L}}}{z-E_r}
  +
  \int_{L_\theta}\rmd E'
  \frac{\ket{\Psi_{E'}^{\mathrm R}}
  \bra{\Psi_{E'}^{\mathrm L}}}{z-E'} .
  \label{eq:csm-resolvent-expansion}
\end{align}
The first sum runs over bound states, the second over exposed resonances, and
$L_\theta$ is the rotated continuum.  This is the operator counterpart of the
contour deformation in Eq.~\eqref{eq:inverse-laplace}.  It also explains why a
finite-basis diagonalization can approximate both pole and continuum
contributions using one matrix spectrum.
